(1) Scaling a representative by $\lambda \neq 0$ scales $s$ and $x$ alike, so $x/s$ is unchanged and $\dsc$ is well defined; the map $N \mapsto x(N)/s(N)$ is continuous on the open, scale-invariant set $\{s \neq 0\}$ of the null cone, so it descends to a continuous map on the quotient. For the lift, $\dsc[X(x)] = x/1 = x$. Conversely, given a non-ideal $[N]$, normalise $s = 1$; the incidence condition \eqref{eq:incidence} then forces $t = \tfrac12 Q(x)$, so the normalised representative is exactly $X(x)$ with $x = \dsc[N]$, whence $X(\dsc[N]) \in [N]$. Both directions are continuous, giving the homeomorphism.

(2) By Theorem \claimref{transport.totality}(2), $\iota[X(x)] = [X(\rho(x))]$, and $\dsc[X(\rho(x))] = \rho(x)$ by (1).

(3) Immediate from Theorem \claimref{transport.totality}(3) and the definition of the ideal points.
